\section{Introduction}
\label{sec:introduction}


Predictive performance in temporal learning depends not only on the learner, but also on the observation protocol under which the data are collected. Standard benchmarks typically hold this protocol fixed and compare learners on the resulting data. This leaves a different question unanswered: what predictive performance could be achieved under an undeployed observation protocol? We refer to this problem as \emph{counterfactual protocol evaluation}. The central question is whether the joint measurement–target distribution under the realised protocol determines the predictive value of an undeployed protocol.


This question is particularly important when the prediction target is defined over a substantially longer time horizon than the observed input. In such settings, the realised protocol may omit temporal information relevant to the target. Examples include predicting whole-night sleep-stage proportions from partially observed sleep records and long-horizon atrial-fibrillation burden from shorter monitoring windows. Weak predictive performance may therefore reflect limitations of the learner, limited sample size, or target-relevant information excluded by the observation protocol. Better models and larger datasets can address the first two limitations, but the third requires changing what is observed.




Several lines of work address related questions. Methods for incomplete or irregular temporal data learn prediction rules from irregularly or incompletely observed training data \citep{che2018grud,rubanova2019latent,kidger2020neural}. Experimental-design and sensor-placement methods, including optimal design for longitudinal data, choose future measurements to optimise a criterion under an assumed or estimated model \citep{chaloner1995bayesian,krause2008near,ji2017optimal}. Bayes-error analyses characterise predictive limits for a given feature--target distribution \citep{fukunaga1987bayes,ishida2023performance}. Most closely related, active feature acquisition performance evaluation studies how to evaluate a new acquisition policy from data collected under a different retrospective acquisition process, under identification assumptions that include sufficient support for the acquisitions required by the target policy \citep{vonkleist2025evaluation,precup2000eligibility,kallus2020double}. In our setting, an alternative protocol may instead require measurements that were never collected under the realised protocol. This leaves a different identification question: whether the joint measurement--target distribution generated by one realised observation protocol determines the predictive value of another protocol whose measurements were never collected.


We develop a theory of counterfactual protocol evaluation centred on the information needed to determine the value of an undeployed protocol. The value of an alternative protocol is identified precisely when it is constant across all latent models that induce the same realised measurement--target law \citep{koopmans1950identification,rothenberg1971identification}. This criterion is value-specific: the value may be identified even when the latent covariance is not, provided that the remaining ambiguity does not change the value being evaluated. Under a Gaussian latent-trajectory model with noisy linear measurements, we show that this condition can fail even with infinite data under the realised protocol, because distinct latent covariances can generate the same joint law of the realised measurements and target while assigning different values to the same alternative. For linear targets, we characterise the unresolved covariance variation through perturbations invisible to the realised law; nonzero directional sensitivity of the alternative's value along such a perturbation certifies local non-identification. Targeted augmentation can remove the value-changing ambiguity and identify the value of a specified alternative without recovering the full latent covariance. A stationary construction gives the sharp minimal failure case in the stationary, standardised, one-observation setting with a mean target, while a permutation construction establishes exact non-identification for nonlinear aggregate targets.


Targeted augmentation may be sufficient when the value of a specified alternative is sought. For an entire candidate family, densely observed calibration trajectories provide a stronger, family-wide route: at the population level, their law identifies the latent covariance used to evaluate all candidates. With only finitely many such trajectories, two questions remain: how accurately candidate values can be distinguished and how those distinctions should guide design. We derive uniform bounds that propagate covariance-estimation error to protocol-value error over the candidate family. The resulting bounds yield selection-regret guarantees and a calibration resolution below which value differences cannot be distinguished reliably. They also link the useful granularity of the candidate class to the amount of calibration data: finer classes can improve the attainable value, but their increasingly small differences support reliable selection only when the data can resolve them. At that resolution, we formulate cost-constrained, target-aware observation design and derive exact rank-one marginal gains for adding measurements under constraints on timing, support, repetition, noise, and cost. Simulations examine calibration error, protocol selection, and design under known data-generating laws. Retrospective analyses of Sleep-EDF and Long-Term AF annotations compare alternative temporal layouts at matched observation budgets \citep{kemp2000analysis,mourtazaev1995age,goldberger2000physiobank,petrutiu2007abrupt}; the pooled analyses show clearer differences between pre-specified dispersed and contiguous layouts than among learned Sleep supports, whose locations and held-out advantages vary across targets and resamples.


Taken together, these results place identification before optimisation in temporal observation design. The information required to evaluate a protocol is not generally the information required to recover the full latent dependence structure; it is the information needed to eliminate ambiguity that changes the protocol values under consideration. Finite calibration then determines the resolution at which those values can be compared, and the granularity of the design problem should be matched to that resolution. This value-directed view links the choice of additional measurements, the amount of calibration data, and the level of detail at which candidate protocols can be compared and optimised.


The remainder of the paper is organised as follows. Section 2 defines temporal aggregate targets, observation protocols, and protocol value. Section 3 develops value-specific identification and targeted augmentation. Section 4 studies finite calibration and protocol comparison, and Section 5 develops target-aware observation design. Section 6 reports the simulations and retrospective analyses. Section 7 reviews related work, and Sections 8 and 9 present the discussion and conclusion.
